\subsection*{99. A 1-Week Truth Test for the D2C Hypothesis}

During a first research week, approach a small set of people who recently handled the relevant AI engineering job, through lawful and appropriate professional channels. Ask about their actual task, deadline, decision criteria, workarounds, private or employer spending, and why available resources failed. Do not start with ``Would you buy our course?''

By the end, the founder should have examples of real recent projects, alternatives, expenses or costs of delay where shareable, and a clearer judgment about whether the buyer needs a course, one-off review, peer community, or nothing external.

If most respondents already have senior technical support and no interest in paying, that is important evidence. It suggests narrowing further, changing payer, or choosing another market rather than creating persuasive landing-page copy.

\subsection*{100. A 2-Week Paid-Proof Test}

Design 1 complete technical experience with a fixed scope and deliverable. Show a public sample task and rubric. Present an actual price and transparent service terms to eligible respondents. Ask for a real paid commitment only when the session can be delivered honestly within the founder's schedule and applicable legal or employment obligations.

Record qualified prospects, offers, purchase decisions, enrollment source, objections, instructional hours, feedback quality, and whether the required work was completed. A buyer who pays but does not engage creates very different evidence from one who finishes and recommends.

A handful of purchases can show possibility and reveal operational issues. It cannot precisely estimate the conversion rate of the whole professional market.

\subsection*{101. The First D2C Cohort Should Be Defined by a Learning Standard}

If the pilot supports a longer offering, make its eligibility criteria educational rather than theatrical: prerequisites, demonstrated willingness to practice, a specific desired capability, and a realistic weekly time commitment. Candidate selection should help teaching work, not imitate famous universities' selectivity without equivalent demand or legitimacy.

The course should have a measurable destination: an engineering artifact, evaluation report, oral defense, and reproducible performance standards where suitable. The cohort should be large enough to provide useful peer learning and commercial contribution but small enough for review quality. The appropriate size is determined by staffing and measured review labor, not a prestige narrative.

\subsection*{102. Do Not Let the First Product Trap Bayes in an Accidental Category}

A program sold as fast job placement may acquire different buyers, sales methods, and institutional expectations from one sold as rigorous adult engineering education. A broad beginner bootcamp may create substantial operational work that cannot easily transition into advanced academic standards. The earliest marketing and product experience establish what the name means.

This does not mean Bayes should avoid career relevance or employability. It means claims should be supported by the capabilities and evidence the institution actually controls. An intensive, well-assessed program can be career-relevant without implying guaranteed salaries or hiring.

The chosen beachhead should support the kind of \textbf{institutional identity the founder wants to compound.}

\subsection*{103. When a College Pilot Should Displace D2C}

Suppose a college provides a named decision-maker, a written and properly authorized funded scope, a feasible semester timetable, adequate student prerequisites, an internal faculty coordinator, reasonable payments and cancellation terms, permission to assess work honestly, and the ability to deliver without conflict with the founder's other obligations.

That is far stronger evidence than 20 hypothetical D2C customers expressing verbal approval. In that case, a bounded college pilot may be the rational primary commercial entry. It can create cash, pedagogical proof, relationships, and opportunities to develop faculty processes, provided the contract economics and academic standards work.

The recommendation is therefore conditional on evidence quality, not loyalty to a channel identity.

\subsection*{104. When Enterprise Might Leapfrog Both}

Likewise, imagine an engineering head offers a paid team assessment on a well-defined AI reliability problem, with appropriate data access, a short approval path, fixed deliverables, a realistic fee, and a scope the founder can deliver sustainably. That opportunity could be a superior first experiment even if the longer-run identity remains educational.

But separate professional consulting revenue from evidence that an education business has product-market fit. A bespoke implementation might fund future curriculum research without proving that people will purchase standardized instruction. Record the assets that transfer and the aspects that will not.

\subsection*{105. The Strategic Value of the First 10 Paying Customers}

In a large institution, 10 customers are a rounding error. In an unknown institute, they are the beginnings of a purchasing mechanism and a reputation. The first 10 suitable paying learners or buyers can reveal actual expectations, where trust is earned, which standards are sustainable, what work customers value, and which aspects of delivery must be improved.

Do not prioritize headcount in the first cohort over customer fit. Ten poor-fit students can consume attention, drop out, and create misleading outcome data. Five unusually appropriate customers whose needs are well understood may produce better learning and stronger ethical proof. But avoid romanticizing tiny groups: the next test must establish whether others beyond the warm founding network will pay.

\subsection*{106. A 2-Axis Bayes Expansion Map}

Bayes may eventually grow along 2 axes: who pays and uses, and what capability the institute delivers:

\begin{table*}[t]
\centering
\small
\begin{tabularx}{\textwidth}{l p{4cm} p{4.5cm} X}
\toprule
\textbf{Expansion Direction} & \textbf{Possible Future Step} & \textbf{What Transfers from D2C Studio} & \textbf{What Must Be Newly Proven} \\
\midrule
Same learner, deeper & Advanced systems-design studio & Learner trust, assessment, operations & New value and repeat willingness to pay \\
Similar learner, adjacent & Quantitative modeling or data systems & Academic standards, faculty processes & Domain demand and specialist instructors \\
Employer sponsorship & Team-based technical assessment & Rubrics and practitioner credibility & Organizational ROI, procurement cycles \\
College partnerships & Faculty-enabled applied lab & Technical content, review standards & Timetable fit, student readiness, approvals \\
Applied research services & Industry experiments and prototypes & Technical methods, relationships & Research economics and long-term delivery \\
Physical learning center & Small in-person studio & Cohort community and pedagogy & Local demand, utilization, fixed facility costs \\
Research-active institute & Labs, research programs, faculty & Culture, talent, curriculum, reputation & Scientific output, grant funding, governance \\
Degree-granting body & Accredited higher education & Academic and institutional assets & Legal authority, accreditation, massive capital \\
\bottomrule
\end{tabularx}
\caption{The 2-axis expansion map for Bayes Institute from initial studio to mature institution.}
\label{tab:ch7-bayes-expansion-map}
\end{table*}

These are \textbf{options, not required steps.} A new institute should not infer that success in an online course automatically authorizes or finances a degree-awarding university.

\subsection*{107. The Expansion Gate I Would Impose on Bayes}

I would not prescribe a universal rupee threshold for adding a second offering. I would require evidence in 5 categories:

\begin{enumerate}
    \item \textbf{Customer evidence:} More than 1 independent group of qualified customers has purchased the core offer at workable realized prices, and outcomes are credible.
    \item \textbf{Economic evidence:} Realistic contribution, cash collection, refund exposure, and normalized instructor costs are documented, with enough liquidity to fulfill every active obligation.
    \item \textbf{Operational evidence:} Teaching and feedback meet standards without permanent emergency founder involvement; a qualified second instructor or reviewer can reproduce at least some results where appropriate.
    \item \textbf{Distribution evidence:} The next customer is not found exclusively through a personal favor; at least 1 channel shows some repeatability, with limits acknowledged.
    \item \textbf{Strategic evidence:} The adjacent offering reuses real assets and has its own observed buyer problem, not merely a shared AI theme.
\end{enumerate}

A low-risk research workshop can be tested before every condition is perfect. A permanent new vertical, hired team, or physical site should face far more stringent requirements.

\subsection*{108. A Conditional First-Year Sequence, Not a Fantasy Roadmap}

\begin{itemize}
    \item \textbf{Phase A: Discover 1 precise job.} Conduct disciplined customer research across a few candidate segments; choose a primary target and a reason to reject weaker alternatives. Publish 1 serious competence artifact and define a bounded paid offer.
    \item \textbf{Phase B: Earn and deliver the first commitments.} Sell a feasible small engagement. Measure actual payment, effort, capability evidence, instructional labor, and refunds. Refuse overpromising and unsafe expansion.
    \item \textbf{Phase C: Repeat the same job.} Attempt a second independent acquisition cycle with clearer positioning and improved operations. Determine whether the result can be reproduced without relying on the same personal network or excessive founder hours.
    \item \textbf{Phase D: Institutionalize the strongest mechanism.} Document curriculum, evaluation, recruiting standards, feedback procedures, sales scripts, finances, and partner obligations. Train another qualified contributor where justified by demand and economics.
    \item \textbf{Phase E: Test exactly 1 meaningful adjacency.} Choose it because buyer evidence and transferable assets make it attractive, not because a new revenue stream sounds impressive. Continue tracking the core rather than allowing the experiment to consume delivery capacity.
\end{itemize}

There is no guarantee that this sequence fits neatly into 12 months. College procurement and a rigorous 6-month program take longer to assess than a workshop. The \textbf{milestones are evidence states, not calendar promises.} The institution's future depends more on the quality of choices than on the theatrical speed with which offerings are announced.
